%=======================================================================
\chapter{Extraction des caractéristiques}

\section{Vue d'ensemble}

L'efficacité d'un système de classification d'images médicales dépend fortement de la qualité
des caractéristiques utilisées pour représenter les informations visuelles contenues dans les
images. Dans le contexte de la classification des tumeurs cérébrales à partir d'images IRM, une
représentation pertinente doit être capable de capturer différents aspects de l'image tels que
l'intensité des pixels, la texture, les variations fréquentielles ainsi que les structures
morphologiques.

Afin d'obtenir une représentation riche et complémentaire des images, ce projet adopte une
approche multi-descripteurs basée sur l'extraction de plusieurs familles de caractéristiques
manuelles (\emph{Handcrafted Features}). Chaque famille permet de décrire une propriété
spécifique de l'image~:

\begin{itemize}[leftmargin=1.4em,itemsep=2pt]
  \item les caractéristiques statistiques décrivent la distribution globale des niveaux de gris~;
  \item les descripteurs de texture (GLCM et LBP) capturent l'organisation spatiale des pixels~;
  \item les caractéristiques fréquentielles extraites par DWT analysent les variations à
        différentes échelles~;
  \item les descripteurs HOG représentent les contours et les orientations dominantes des
        gradients.
\end{itemize}

Cette diversité permet d'exploiter simultanément plusieurs sources d'information et d'étudier
leur complémentarité dans le processus de classification.

La Figure~\ref{fig:extraction-pipeline} illustre l'architecture générale de l'extraction des
caractéristiques utilisée dans ce travail.

% \begin{notebox}[Pipeline d'extraction~:]
% Image prétraitée $\rightarrow$ extraction parallèle des branches A (Stats), B (GLCM+LBP),
% C (DWT) et D (HOG) $\rightarrow$ fusion des vecteurs $\rightarrow$ classifieur ML.
% \end{notebox}

\begin{figure}[htbp]
  \centering
  \begin{tikzpicture}[
    x=1cm, y=1cm,
    input/.style={rectangle, rounded corners=3pt, draw=accent, fill=headbg,
                  minimum width=2.8cm, minimum height=0.95cm, align=center,
                  font=\small\bfseries},
    branch/.style={rectangle, rounded corners=3pt, draw=accent, fill=white,
                   minimum width=2.15cm, minimum height=1.25cm, align=center,
                   font=\footnotesize},
    merge/.style={rectangle, rounded corners=3pt, draw=accent2, fill=headbg,
                  minimum width=2.8cm, minimum height=0.95cm, align=center,
                  font=\small\bfseries},
    clfbox/.style={rectangle, rounded corners=3pt, draw=accent, fill=white,
                   minimum width=2.4cm, minimum height=0.9cm, align=center,
                   font=\footnotesize},
    arr/.style={-{Stealth[length=2mm]}, thick, accent!85}
  ]
    % --- nodes ---
    \node[input] (img) at (0,0) {Image\\prétraitée};

    \node[branch] (a) at (-4.5,-2.1) {A --- Stats\\[1pt]
      {\scriptsize 11 dim.}};
    \node[branch] (b) at (-1.5,-2.1) {B --- Texture\\[1pt]
      {\scriptsize GLCM+LBP}\\{\scriptsize 20 dim.}};
    \node[branch] (c) at (1.5,-2.1) {C --- DWT\\[1pt]
      {\scriptsize 16 dim.}};
    \node[branch] (d) at (4.5,-2.1) {D --- HOG\\[1pt]
      {\scriptsize 1944 dim.}};

    \node[merge] (fuse) at (0,-3.75) {Fusion\\[1pt]
      {\scriptsize Full = 1991 dim.}};
    \node[clfbox] (clf) at (0,-5.05) {Classifieur\\[1pt]
      {\scriptsize SVM, KNN, \ldots}};

    % --- top fan-out ---
    \coordinate (topbus) at (0,-0.75);
    \draw[arr] (img.south) -- (topbus);
    \draw[arr] (topbus) -| (a.north);
    \draw[arr] (topbus) -| (b.north);
    \draw[arr] (topbus) -| (c.north);
    \draw[arr] (topbus) -| (d.north);

    % --- bottom fan-in ---
    \coordinate (botbus) at (0,-3.05);
    \draw[arr] (a.south) -- (a.south |- botbus);
    \draw[arr] (b.south) -- (b.south |- botbus);
    \draw[arr] (c.south) -- (c.south |- botbus);
    \draw[arr] (d.south) -- (d.south |- botbus);
    \draw[arr] (-4.5,-3.05) -- (4.5,-3.05);
    \draw[arr] (botbus) -- (fuse.north);

    \draw[arr] (fuse.south) -- (clf.north);
  \end{tikzpicture}
  \caption{Architecture générale du pipeline multi-descripteurs.}
  \label{fig:extraction-pipeline}
\end{figure}

% -----------------------------------------------------------------------
\section{Histogram of Oriented Gradients (HOG)}

\subsection{Principe}

Le descripteur \emph{Histogram of Oriented Gradients} (HOG) est une méthode largement utilisée
pour représenter les formes et les contours présents dans une image.

Le principe consiste à calculer les gradients locaux puis à construire des histogrammes
représentant la distribution des orientations dominantes.

Les régions tumorales présentent souvent des structures morphologiques spécifiques qui peuvent
être efficacement capturées par cette représentation.

\subsection{Hyperparamètres étudiés}

\paragraph{Orientations}
Nombre de directions utilisées pour construire l'histogramme des gradients. Une valeur élevée
permet une représentation plus détaillée des contours.

\paragraph{Pixels per Cell}
Détermine la taille des cellules élémentaires utilisées pour calculer les histogrammes locaux.
Ce paramètre influence directement le niveau de détail capturé.

\paragraph{Cells per Block}
Les cellules sont regroupées en blocs afin d'effectuer une normalisation locale. Cette étape
améliore la robustesse du descripteur face aux variations de contraste.

La configuration optimale retenue est~: \hogParams{} (score composite~: \hogScore).

Le vecteur HOG obtenu fournit une représentation détaillée des structures géométriques
présentes dans les images IRM et complète efficacement les informations statistiques,
texturales et fréquentielles extraites par les autres descripteurs.

\begin{figure}[htbp]
  \centering
  \includegraphics[width=\textwidth]{hog_schema}
  \caption{Schéma illustratif du calcul HOG~: image d'entrée, visualisation des gradients
    orientés et vecteur de caractéristiques résultant.}
  \label{fig:hog-schema}
\end{figure}

% -----------------------------------------------------------------------
\section{Caractéristiques statistiques (A)}

Les caractéristiques statistiques constituent la forme la plus simple de représentation d'une
image. Elles permettent de résumer la distribution globale des intensités des pixels à travers
plusieurs indicateurs numériques.

Bien que relativement simples, ces mesures fournissent des informations importantes sur la
luminosité, la dispersion et la forme de la distribution des niveaux de gris. Elles sont
souvent utilisées comme descripteurs de base dans les systèmes de classification médicale.

Pour chaque image, les caractéristiques suivantes sont extraites~:

\paragraph{Moyenne (Mean)}
La moyenne représente l'intensité moyenne des pixels de l'image~:
\begin{equation}
  \mu = \frac{1}{N}\sum_{i=1}^{N} x_i
\end{equation}
où $x_i$ représente l'intensité du pixel $i$ et $N$ le nombre total de pixels. Cette mesure
fournit une indication globale du niveau de luminosité de l'image.

\paragraph{Variance}
La variance mesure la dispersion des intensités autour de la moyenne~:
\begin{equation}
  \sigma^2=\frac{1}{N}\sum_{i=1}^{N}(x_i-\mu)^2
\end{equation}
Une variance élevée indique une forte hétérogénéité des niveaux de gris.

\paragraph{Écart-type (Standard Deviation)}
L'écart-type correspond à la racine carrée de la variance~:
\begin{equation}
  \sigma=\sqrt{\sigma^2}
\end{equation}
Il constitue un indicateur direct de la variabilité des intensités.

\paragraph{Entropie}
L'entropie mesure la diversité de la distribution d'intensités (calculée sur un histogramme
normalisé à 64 bins).

\paragraph{Asymétrie (Skewness)}
L'asymétrie mesure le degré de dissymétrie de la distribution des niveaux de gris. Une valeur
positive indique une distribution étirée vers les fortes intensités tandis qu'une valeur
négative traduit une concentration vers les faibles intensités.

\paragraph{Kurtosis (Aplatissement)}
Le kurtosis mesure le degré de concentration des valeurs autour de la moyenne. Cette
caractéristique permet de distinguer les distributions présentant des pics prononcés de celles
plus étalées.

Le vecteur implémenté (\texttt{feat\_stats}) regroupe exactement \dimStats{} descripteurs~:
moyenne, écart-type, variance, entropie, asymétrie, kurtosis et les percentiles
10e, 25e, 50e (médiane), 75e et 90e.

L'ensemble de ces mesures constitue le groupe de caractéristiques \textbf{A --- Statistical
Features} utilisé dans la suite du projet.

% -----------------------------------------------------------------------
\section{Gray Level Co-occurrence Matrix (GLCM)}

\subsection{Principe}

La matrice de cooccurrence des niveaux de gris (\emph{Gray Level Co-occurrence Matrix} --- GLCM)
est l'une des méthodes les plus utilisées pour l'analyse de texture.

Introduite par Haralick, cette technique repose sur l'étude des relations spatiales entre les
pixels voisins. La matrice GLCM comptabilise la fréquence d'apparition de paires de niveaux de
gris séparées par une distance et une direction données.

Ainsi, au lieu d'étudier chaque pixel individuellement, la méthode analyse la structure
spatiale de l'image, ce qui permet de caractériser efficacement les textures présentes dans les
régions tumorales.

\subsection{Hyperparamètres étudiés}

Dans le cadre de ce projet, plusieurs configurations de GLCM ont été évaluées lors de la phase
d'optimisation (Phase~1).

\paragraph{Distances}
La distance définit l'écart entre les pixels considérés lors du calcul de la matrice. Les
distances testées permettent de capturer des textures locales ou plus globales selon l'échelle
étudiée.

\paragraph{Angles}
Les directions utilisées influencent la manière dont les relations spatiales sont observées.
Les orientations classiques comprennent~:
\begin{equation}
  0^\circ,\ 45^\circ,\ 90^\circ,\ 135^\circ
\end{equation}
Ces différentes directions permettent de prendre en compte l'anisotropie potentielle des
textures.

\paragraph{Levels}
Le paramètre \emph{levels} définit le nombre de niveaux de gris utilisés lors de la
quantification de l'image. Une quantification appropriée permet de réduire le bruit tout en
conservant les structures texturales importantes.

\subsection{Caractéristiques extraites}

À partir de chaque matrice GLCM, plusieurs indicateurs statistiques sont calculés.

\paragraph{Contrast}
Mesure les variations locales des niveaux de gris. Des valeurs élevées traduisent généralement
des textures fortement contrastées.

\paragraph{Correlation}
Évalue la dépendance linéaire entre les pixels voisins. Cette mesure permet d'apprécier
l'organisation spatiale de la texture.

\paragraph{Energy}
Mesure l'uniformité de la texture (somme des carrés des probabilités de cooccurrence). Une
énergie élevée indique une texture homogène.

\paragraph{Homogeneity}
Mesure la proximité des valeurs de la matrice par rapport à sa diagonale principale. Une
texture homogène présente généralement des valeurs élevées de cette caractéristique.

\paragraph{Dissimilarity et ASM}
La \emph{dissimilarity} quantifie la variation locale linéaire des niveaux de gris, tandis que
l'\emph{ASM} (\emph{Angular Second Moment}) est une mesure complémentaire d'uniformité
texturale. Au total, six descripteurs GLCM sont conservés (moyennés sur toutes les paires
distance/angle)~: contrast, dissimilarity, homogeneity, energy, correlation et ASM.

La configuration optimale retenue est~: \glcmParams{} (score composite~: \glcmScore).

\begin{figure}[htbp]
  \centering
  \includegraphics[width=\textwidth]{glcm_best_params}
  \caption{Illustration de l'extraction GLCM~: image d'entrée et valeurs des descripteurs de
    texture calculés (meilleurs paramètres Phase~1).}
  \label{fig:glcm-best}
\end{figure}

% -----------------------------------------------------------------------
\section{Local Binary Patterns (LBP)}

\subsection{Principe}

Le descripteur \emph{Local Binary Patterns} (LBP) est une méthode efficace d'analyse de texture
locale.

Son principe consiste à comparer chaque pixel central avec ses voisins. Si l'intensité d'un
voisin est supérieure ou égale à celle du pixel central, un bit de valeur~1 est attribué~;
sinon un bit de valeur~0 est attribué.

La concaténation de ces bits forme un code binaire décrivant la structure locale du voisinage.
L'histogramme des codes LBP obtenus constitue ensuite le vecteur de caractéristiques de
l'image.

Cette méthode est particulièrement robuste aux variations d'éclairage et largement utilisée
pour la reconnaissance de textures.

\subsection{Hyperparamètres étudiés}

\paragraph{$P$ (Nombre de voisins)}
Le paramètre $P$ représente le nombre de points de voisinage utilisés pour construire le motif
binaire.

\paragraph{$R$ (Rayon)}
Le rayon $R$ définit la distance entre le pixel central et les pixels voisins. Des rayons plus
élevés permettent de capturer des structures plus larges.

\paragraph{Method}
Plusieurs variantes de LBP ont été testées dans la grille Phase~1~: \emph{uniform},
\emph{nri\_uniform} et \emph{ror} (\emph{rotation invariant}). La variante \emph{uniform}
retenue produit un histogramme de $P+2$ bins (soit \dimLBP{} descripteurs pour $P=12$).

La configuration optimale retenue est~: \lbpParams{} (score composite~: \lbpScore).

\begin{figure}[htbp]
  \centering
  \includegraphics[width=\textwidth]{lbp_best_params}
  \caption{Illustration de l'extraction LBP~: image d'entrée, carte LBP et histogramme des
    motifs (meilleurs paramètres Phase~1).}
  \label{fig:lbp-best}
\end{figure}

% -----------------------------------------------------------------------
\section{Discrete Wavelet Transform (DWT)}

\subsection{Principe}

La Transformée en Ondelettes Discrètes (DWT) est une méthode d'analyse fréquentielle
permettant de représenter une image simultanément dans les domaines spatial et fréquentiel.

Contrairement à la transformée de Fourier, la DWT conserve une localisation spatiale des
informations, ce qui la rend particulièrement adaptée à l'analyse des images médicales.

Lors de la décomposition, l'image est séparée en plusieurs sous-bandes~:

\begin{itemize}[leftmargin=1.4em,itemsep=2pt]
  \item approximation (LL)~;
  \item détails horizontaux (LH)~;
  \item détails verticaux (HL)~;
  \item détails diagonaux (HH).
\end{itemize}

Ces sous-bandes permettent de capturer des informations texturales à différentes échelles.

\subsection{Familles d'ondelettes}

La grille Phase~1 teste les ondelettes suivantes~: \texttt{haar}, \texttt{db1}, \texttt{db2},
\texttt{db4}, \texttt{sym2}, \texttt{coif1} et \texttt{bior1.3}. Pour chaque sous-bande, quatre
statistiques sont calculées~: moyenne, écart-type, énergie et entropie.

\subsection{Niveaux de décomposition}

Le niveau de décomposition contrôle la profondeur de l'analyse multi-résolution. Des niveaux
plus élevés permettent de capturer des structures globales tandis que des niveaux faibles
privilégient les détails locaux.

Les statistiques calculées sur les coefficients d'ondelettes constituent les caractéristiques
DWT utilisées pour la classification.

La configuration optimale retenue est~: \dwtParams{} (score composite~: \dwtScore).

\begin{figure}[htbp]
  \centering
  \includegraphics[width=\textwidth]{dwt_best_params}
  \caption{Illustration de l'extraction DWT~: image d'entrée, sous-bande d'approximation et
    vecteur de caractéristiques (meilleurs paramètres Phase~1).}
  \label{fig:dwt-best}
\end{figure}

% -----------------------------------------------------------------------
\section{Synthèse des dimensions par groupe de caractéristiques}

Le Tableau~\ref{tab:feature-dims} récapitule la dimension de chaque groupe de
caractéristiques (A, B, C, D) avant l'étape de nettoyage, pour la configuration optimale retenue
après la Phase~1.

\begin{table}[htbp]
  \centering
  \caption{Dimensions des groupes de caractéristiques avant nettoyage.}
  \label{tab:feature-dims}
  \begin{tabular}{@{}clcr@{}}
    \toprule
    \thd{Groupe} & \thd{Contenu} & \thd{Détail} & \thd{Dim.} \\
    \midrule
    A & Statistiques & mean, std, var, entropy, skew, kurtosis, p10--p90 & \dimStats \\
    B & Texture & GLCM (\dimGLCM) + LBP (\dimLBP) & \dimTexture \\
    C & DWT & Statistiques sur sous-bandes (niveau~1) & \dimDWT \\
    D & HOG & Histogrammes de gradients orientés & \dimHOG \\
    \midrule
    \textbf{Full} & \textbf{Fusion A+B+C+D} & & \textbf{\dimFull} \\
    \bottomrule
  \end{tabular}
\end{table}

La branche HOG constitue de loin le descripteur le plus volumineux, ce qui justifie
l'importance des étapes de nettoyage et de réduction de redondance décrites au
Chapitre~\ref{ch:ensembles}.
